\begin{answers}{ch:agent}{AI エージェントを作る}
\item 14 個です。
緑は、モデルが書く 2 つの発話です。
呼び出しの発話は、\code{<|tool|>} が 1 個、「\code{48*55}」が 5 個、\code{<|end|>} が 1 個で、合わせて 7 個です。
最終応答の発話も、\code{<|assistant|>} が 1 個、「\code{2640です。}」が 5 個、\code{<|end|>} が 1 個で、7 個です。
利用者の発話と、電卓が書く \code{<|result|>2640<|end|>} は損失に入りません。
\code{learn_role=True} を付けなければ、役割トークンの 2 個が外れて 12 個になります。
\item 段落 1 と段落 3 が同点になり、一番近い段落を 1 つに決められません。
ベクトルの長さは、段落 1 と段落 3 が $\sqrt2$、質問が $\sqrt3$ です。
どちらの段落も質問と 2 つの組を共有するので、内積はどちらも $\frac{1\times1+1\times1}{\sqrt2\,\sqrt3}=\frac{2}{\sqrt6}\approx0.82$ です。
出現回数だけでは、「共有する組がいくつあるか」しか見ていないからです。
idf を掛けると、2 つの段落に出る「の日」（0.41）より、1 つの段落にしか出ない「の夜」（1.10）の一致が重く数えられ、類似度は 0.46 と 0.94 に分かれます。
idf の役割は、珍しい組の一致を、ありふれた組の一致より重く数えて、内容を見分ける手がかりにすることです。
\item 解けるようになりました（実測）。
\code{TEMPLATES} に「\code{a×b×cは？}」と「\code{a×bにcを掛けるといくつですか？}」を足し、式を \code{a*b*c} と書く手本を加えて、4 種類を 1,500 件ずつで SFT しました。
「\code{3×4×5は？}」には \code{<|tool|>3*4*5<|end|>} と書き、「\code{60です。}」と答えました。
掛けて掛ける問題 100 問の正答率は、\file{ch23_tool_sft.pt} の 1\% から 97\% に上がりました。
ほかの 3 種類も、加減算 79\%、掛け算 100\%、掛けて足す問題 100\% で、崩れていません。
間違えた 3 問のうち 2 問は、8〜9 桁の結果を答えに写すときの間違い（結果の誤読）でした。
手本に無い「\code{2×3×4×5は？}」も \code{2*3*4*5} と書けましたが、「\code{7の2乗は？}」には \code{72*2} と書きました。
手本を足した形は直りますが、手本の範囲の外で崩れることは変わりません。
\item 3 通りのどの言い回しでも、電卓を呼ばずに 8 割強を解くようになりました（実測）。
\code{make_tool_example} の加減算の質問を、\code{QUESTION_TEMPLATES} の 3 通りから選ぶように変え、同じ設定で SFT しました。
\tabref{tab:agent-phrasing}の 3 行は、電卓を呼んだ割合がどれも 0.0\%、正答率が上から 85.3\%、84.0\%、82.0\% になりました（元は 80.7\%、4.7\%、5.3\%）。
掛け算は 96\%、掛けて足す問題は 92\% で、どちらも毎回電卓を呼んでおり、使い分けは保たれています。
ただし、直ったのは手本に入れた 3 通りだけです。
手本に無い「\code{aプラスbは？}」では、150 問すべてで電卓を呼んで式を \code{*} で書き、正答率は 0\% でした。
また、3 通りで残った誤答（450 問中 73 問）は、すべて暗算の間違いでした。これは言い回しの手本を足しても減りません。
電卓に任せるべきかどうかを、\chapref{ch:agentic-rl}では結果の採点から学ばせます。
\item \code{eval} なら、台本が渡した \code{__import__('os').system('ls')} がそのまま実行されます。
\code{ls} が本当に動いてファイルの一覧が画面に出て、文脈には命令の成功を表す 0 が \code{<|result|>0<|end|>} として足されました（電卓を \code{eval} に替えて実測）。
\code{ls} が \code{rm -rf ~} なら、ファイルが消えます。
\code{calculator} には、関数呼び出し、変数名、\code{9**9**9}、小数、\code{True}、文字列、全角の数字など 56 通りの入力を試しました。
コードが実行される入力も、例外が外に出る入力も、見つかりませんでした。
通るのは整数と \code{+}・\code{-}・\code{*} の節点だけで、式も 40 字までなので、結果は多くても 40 桁ほどの整数です。
エラー文の「使える記号は 0-9 + - * のみ」より少し広く、括弧、空白、\code{1_000}、\code{0x10}（16 と読まれる）も通りますが、整数の計算にしかならないので害はありません。
見つからなかったことは、すり抜けが無いことの証明ではありません。
それでも、「危ないものを探して禁じる」のではなく「許すものだけを通す」作りなので、整数と 3 つの演算でできていない式は、中身が何であれ弾かれます。
\end{answers}
